\documentclass[slidestop,compress,mathserif]{beamer}
\usepackage{fontspec,xunicode,xltxtra,beamerthemesplit}

%以下是各种演示主题，定义幻灯片中的所有细节
%\usetheme{default}
%\usetheme{Berlin}%这个主题比较好
%\usetheme{Pittsburgh}
%\usetheme{Rochester}
\usetheme{Berkeley}%这种演示主题比较好
%\usetheme{Goettingen}
%\usetheme{Hannover}
%\usetheme{Marburg}
%\usetheme{PaloAlto}%这种演示主题比较好
%\usetheme{Antibes}
%\usetheme{Darmstadt}
%\usetheme{JuanLesPins}
%\usetheme{Montpellier}
%\usetheme{Singapore}
%\usetheme{Boadilla}
%\usetheme{Madrid}
%\usetheme{AnnArbor}
%\usetheme{CambridgeUS}
%\usetheme{Copenhagen}
%\usetheme{Warsaw}

%以下是各种外部主题，确定幻灯的显示式样
%\useoutertheme{infolines}
%\useoutertheme{miniframes}
%\useoutertheme[height=0.1\textwidth,width=0.15\textwidth,hideothersubsections]{sidebar}
%\useoutertheme{smoothbars}
%\useoutertheme{split}
%\useoutertheme{shadow}
%\useoutertheme{tree}
%\useoutertheme{smoothtree}
%\useoutertheme[height=0.5\textwidth]{sidebar}

%以下是各种内部主题
%\useinnertheme{default}
%\useinnertheme{circles}
%\useinnertheme{rectangles}
\useinnertheme[shadow]{rounded}

%以下是各种颜色主题
%\usecolortheme{default}
%\usecolortheme{albatross}
%\usecolortheme{beaver}
%\usecolortheme{beetle}
%\usecolortheme{crane}
%\usecolortheme{dolphin}
%\usecolortheme{dove}
%\usecolortheme{fly}
%\usecolortheme{lily}
%\usecolortheme{orchid}
%\usecolortheme{rose}
%\usecolortheme{seagull}
\usecolortheme{seahorse}
%\usecolortheme{sidebartab}
%\usecolortheme{structure}
%\usecolortheme{whale}
%\usecolortheme{wolverine}

%以下是各种字体主题
%\usefonttheme{default}
%\usefonttheme[onlymath]{serif}
%\usefonttheme{structurebold}
%\usefonttheme{structureitalicserif}
%\usefonttheme{structuresmallcapsserif}

\setsansfont[Mapping=tex-text, BoldFont={Microsoft YaHei}]{Microsoft YaHei}

% 中文环境自动换行
\XeTeXlinebreaklocale "zh"  % 表示用中文的断行
\XeTeXlinebreakskip = 0pt plus 1pt % 多一点调整的空间

% 中文环境修正导航栏
\makeatletter
\setbeamertemplate{blocks}[rounded][shadow=true] 
\def\beamer@linkspace#1{
  \begin{pgfpicture}{0pt}{-1.5pt}{#1}{5.5pt}
    \pgfsetfillopacity{0}
    \pgftext[x=0pt,y=-1.5pt]{.}
    \pgftext[x=#1,y=5.5pt]{.}
  \end{pgfpicture}}
\makeatother

% 超链接高亮显示
\hypersetup{CJKbookmarks=true,
colorlinks=true,
citecolor=blue,
linkcolor=blue,
urlcolor=blue,
bookmarksopen=true,
breaklinks=true
}

% 幻灯片切换方式
%\transblindshorizontal 	% 水平百叶窗
%\transblindsvertical 		% 垂直百叶窗
%\transboxin				% 盒状收缩
%\transboxout				% 盒状展开
%\transdissolve				% 溶解
%\transglitter
%\transsplithorizontalin	% 上下向中央收缩
%\transsplitverticalin		% 垂直向中央收缩
%\transsplithorizontalout	% 上下向中央展开
%\transsplitverticalout		% 垂直向中央展开
%\transwipe					% 从下抽出


\title{2002年图灵奖~\\~公钥密码学~\\~RSA加密算法}
\author{刘正~~徐小奇}
\date{\today}
%\institute{同济大学电信学院}

\logo{\color{blue!50}\scalebox{2}{{\includegraphics[height=0.8cm]{logo.jpg}\vspace{220pt}}}}

% 设定frametitle居中
\makeatletter 
\long\def\beamer@@frametitle[#1]#2{% 
  \beamer@ifempty{#2}{}{% 
    \gdef\insertframetitle{\centering{#2\ifnum\beamer@autobreakcount>0\relax{}\space\usebeamertemplate*{frametitle continuation}\fi}}% 
  \gdef\beamer@frametitle{#2}% 
  \gdef\beamer@shortframetitle{#1}% 
}% 
} 
\makeatother


\begin{document}

\frame{\titlepage}

\section{RSA加密算法}

\begin{frame}
  \frametitle{加密的历史}
  
  \begin{block}<+->{1976年以前，所有的加密方法都是同一种模式：}
    \begin{itemize}[<+->]
        \item 甲方选择某一种加密规则，对信息进行加密；
        \item 乙方使用同一种规则，对信息进行解密。
    \end{itemize}
  \end{block}

\end{frame}

\begin{frame}
  \frametitle{加密的历史}
  
  \begin{block}{1976年以前，所有的加密方法都是同一种模式：}
    \begin{itemize}
        \item 甲方选择某一种加密规则，对信息进行加密；
        \item 乙方使用同一种规则，对信息进行解密。
    \end{itemize}
  \end{block}
  ~\\[0.6cm]

~ ~ ~ ~由于加密和解密使用同样规则（简称"密钥"），这被称为"对称加密算法"（Symmetric-key algorithm）。

~ ~ ~ ~这种加密模式有一个最大弱点：甲方必须把加密规则告诉乙方，否则无法解密。保存和传递密钥，就成了最头疼的问题。
\end{frame}


\begin{frame}
  \frametitle{RSA加密算法}
~ ~ ~ ~ RSA加密算法是一种非对称加密算法。在公开密钥加密和电子商业中RSA被广泛使用。

~ ~ ~ ~ RSA是1977年由罗纳德·李维斯特（Ron Rivest）、阿迪·萨莫尔（Adi Shamir）和伦纳德·阿德曼（Leonard Adleman）一起提出的。当时他们三人都在麻省理工学院工作。RSA就是他们三人姓氏开头字母拼在一起组成的。

~ ~ ~ ~ 1973年，在英国政府通讯总部工作的数学家克利福德·柯克斯（Clifford Cocks）在一个内部文件中提出了一个相同的算法，但他的发现被列入机密，一直到1997年才被发表。

\end{frame}

\section{获奖者生平}
\begin{frame}
  \frametitle{Ronald L. Rivest}

\end{frame}

\begin{frame}
  \frametitle{Adi Shamir}

\end{frame}

\begin{frame}
  \frametitle{Leonard M. Adleman}

\end{frame}

\begin{frame}
  \frametitle{Clifford Cocks}

\end{frame}




\section{八卦环节}









\begin{frame}
\frametitle{我是中文} 
%\framesubtitle{\centerline{subtext}} 
\animate<3-6>% 自动逐步显示
\begin{itemize}[<+->]
\item one
\item two
\item three
\item four
\item five
\item six
\end{itemize}

刘正 徐小奇 e
\end{frame}


















\end{document}
